\documentclass[11pt]{article}
\usepackage[margin=1in]{geometry}
\usepackage{amsmath,amssymb,microtype}
\usepackage[hidelinks]{hyperref}

\newcommand{\R}{\mathbb R}
\newcommand{\Szero}{\mathcal S_0^n}
\newcommand{\Sb}{\mathcal S_b^n}

\title{Both star-body valuation questions in arXiv:1611.03345\\
are answered by arXiv:1709.08959}
\author{Literature-status note}
\date{17 August 2026}

\begin{document}
\maketitle

\begin{center}
\fbox{\parbox{0.91\textwidth}{\textbf{Status.}
Explicit affirmative later-paper answers to Questions 7.5 and 7.6.  The later
results are stated on the larger class of bounded Borel star sets and therefore
imply the requested conclusions for star bodies.  This note records a
literature identification, not a new theorem.}}
\end{center}

\section{The two source questions}

Let $\Szero$ be the star bodies in $\R^n$, identified with their continuous
radial functions on $S^{n-1}$, and let $\Sb$ denote bounded Borel star sets.
Tradacete and Villanueva prove in arXiv:1611.03345v1 that a radial continuous
valuation extends from $\Szero$ to $\Sb$ and has an integral representation on
simple radial functions \cite{TV1}.  They then ask on PDF page 23:

\begin{quote}
\textbf{Question 7.5.} Is every radial continuous valuation on the star bodies
of $S^{n-1}$ uniformly continuous on bounded sets?

\textbf{Question 7.6.} Is the integral representation in Theorem 1.2 valid for
every star body?
\end{quote}

The source notes that an affirmative answer to Question 7.5 would imply one to
Question 7.6.

\section{The explicit later answers}

The same authors' follow-up arXiv:1709.08959v1 is titled
\emph{Continuity and representation of valuations on star bodies}
\cite{TV2}.  Its abstract already announces both conclusions.

First, Theorem 3.8 on PDF page 11 states that every radial continuous
valuation $V:\Sb\to\R$ is uniformly continuous on bounded sets: for every
$\lambda,\varepsilon>0$ there is $\delta>0$ such that
\[
 \|f\|_\infty,\|g\|_\infty\leq\lambda,
 \qquad \|f-g\|_\infty\leq\delta
 \quad\Longrightarrow\quad
 |\widetilde V(f)-\widetilde V(g)|\leq\varepsilon.
\]
Restricting this theorem to continuous radial functions gives the affirmative
answer to Question 7.5.

Second, Theorem 1.1 on PDF page 2 characterizes all radial continuous
valuations.  It says that such a valuation is exactly one for which there are
a finite Borel measure $\mu$ on $S^{n-1}$ and a strong-Carath\'eodory kernel
$K:\R_+\times S^{n-1}\to\R$, locally dominated by an $L^1(\mu)$ function,
such that
\[
 V(L)=\int_{S^{n-1}}K(\rho_L(t),t)\,d\mu(t)
\]
for every star set $L$ with bounded Borel radial function.  In particular this
holds for every star body, affirmatively answering Question 7.6 in a stronger
form.

Thus the exact disposition is
\[
 \boxed{\text{Questions 7.5 and 7.6: both affirmative by
 arXiv:1709.08959, Theorems 3.8 and 1.1.}}
\]

\begin{thebibliography}{9}
\bibitem{TV1}
P.~Tradacete and I.~Villanueva,
\emph{Radial continuous valuations on star bodies and star sets},
arXiv:1611.03345v1 (2016).

\bibitem{TV2}
P.~Tradacete and I.~Villanueva,
\emph{Continuity and representation of valuations on star bodies},
arXiv:1709.08959v1 (2017).
\end{thebibliography}

\end{document}

